% ============================================================
% Decision Gate Metrics — proposal_18E/decision_gate.tex
% Reviewer concern: RF-5a (Decision Gate Metrics TODO)
% Verified by: AC-5
%
% Framed as SCALING-BEHAVIOR CLAIMS (not fixed-number commitments).
% Thresholds are discovered empirically via the calibration studies
% in Methods §Phase 1 (AC-5b) and §Phase 2 (AC-5c).
% ============================================================

% [RF-5a] Decision Gate Metrics as scaling-behavior claims with empirically-discovered thresholds — verified by: AC-5

\paragraph{AI advantage as scaling behavior.} Our primary AI advantage metric is the \textbf{monotonic improvement} of provenance completeness, reproducibility, and span-level alignment F1 as the software engineering corpus underlying PA-AKG scales from the Phase 1 small-seed calibration ($\sim$5--10 papers, Graph-Based RAG snowball seed) to the Phase 2 expanded-corpus calibration ($\sim$50+ papers). The comparison of Phase 1 against Phase 2 thresholds provides the direct empirical evidence that additional curated training and retrieval data translates into measurable improvements in the trustworthiness of PA-AKG's outputs. This directly addresses the CFP's call for ``scaling behavior which shows increasing performance as additional data, computing, and/or other resources are applied.''

\paragraph{Empirically-discovered thresholds.} Rather than pre-committing to first-pass numeric targets, the Phase I go/no-go thresholds for provenance completeness, reproducibility rate, and span-level alignment F1 are \emph{discovered} through the calibration studies in \S\ref{sec:methods}. Phase 1 calibration produces an initial threshold set published at the end of Month 3; Phase 2 re-calibration re-validates and refines these against the expanded corpus. The values reported at the \textbf{Month-6 Decision Gate} are the Phase 2 calibrated thresholds, published with the gate report so that Phase II Genesis reviewers can assess whether PA-AKG has cleared its own empirically-grounded bar.

\paragraph{Iterate-if-low rule.} If the Phase 1 or Phase 2 calibration study produces thresholds below a defensible trustworthiness floor---operationally, below the point at which human HPC reviewers consistently judge recommendations as auditable---the team will iterate the PA-AKG implementation (schema, extraction pipeline, or provenance inheritance rules) and re-calibrate before the Month-6 gate, rather than accepting a weak result as the project's claimed AI advantage. This rule protects the integrity of the scaling-behavior claim: improvements must come from additional curated data, not from lowering the bar.

\paragraph{Supporting user-study signal.} As an early user-feedback signal supporting the scaling-behavior claim, controlled studies with HPC software developers will assess recommendation auditability, with a target of $\geq$70\% of developers rating PA-AKG recommendations as auditable, compared against baseline LLM recommendations. This is reported as a supporting metric, not as the primary trustworthiness number.

\paragraph{Provenance-driven acceleration.} PA-AKG's persistent knowledge layer creates a direct acceleration advantage: AI agents resolve queries from KG-stored validated knowledge without re-executing the full retrieval and inference pipeline. By Month 6, we target $\geq$30\% reduction in mean response latency for transformation queries whose answers are already encoded as validated provenance in the KG, compared to baseline RAG systems that treat every query independently. This metric directly captures the compounding benefit of PA-AKG's knowledge state: as the KG grows, the fraction of queries answered from validated provenance increases, reducing the lag between new research findings and actionable code optimizations across DOE R\&D workflows.
